\documentclass[12pt]{article}

\usepackage{preamble}

% ---------------------------------------------------------------------------
% Document-specific macros
% ---------------------------------------------------------------------------
% Inline author notes, e.g. \anote{check this bound}
\newcommand{\anote}[1]{\textcolor{red}{\small\textbf{(Ashwin:} #1\textbf{)}}}
% Generic placeholder for unfinished material
\newcommand{\todo}[1]{\textcolor{\mysubcolor}{\small\textbf{[TODO:} #1\textbf{]}}}

% Algebras, spaces and maps
\newcommand{\alg}{\mathcal{A}}             % domain algebra
\newcommand{\algB}{\mathcal{B}}            % target algebra
\newcommand{\hilb}{\mathcal{H}}            % input Hilbert space
\newcommand{\dil}{\mathcal{K}}             % dilation space
\newcommand{\bdd}[1]{\mathcal{B}\br{#1}}   % bounded operators B(H)
\newcommand{\cp}{\mathrm{CP}}              % completely positive maps
\newcommand{\matr}[2]{\mathrm{M}_{#1}\br{#2}} % n x n matrices over an algebra
\newcommand{\algtensor}{\odot}             % algebraic tensor product
\newcommand{\op}{^{\mathrm{op}}}           % opposite algebra

% Inner products (linear in the first argument)
\newcommand{\ip}[2]{\left\langle #1 , #2 \right\rangle}
\newcommand{\ipphi}[2]{\left\langle #1 , #2 \right\rangle_{\varphi}}
\newcommand{\rankone}[2]{\theta_{#1, #2}}  % w |-> <w, #2> #1

% Standard form of a finite-dimensional algebra
\newcommand{\Ltwo}{L^2\br{\alg}}           % L^2(A, tau)
\newcommand{\vect}[1]{\widehat{#1}}        % a |-> \hat{a} in L^2(A)
\newcommand{\Ltn}{\hilb \tensor \Ltwo}     % H (x) L^2(A)
\newcommand{\Fphi}{F_{\varphi}}            % operator implementing <.,.>_phi
\newcommand{\Choi}{C_{\varphi}}            % Choi matrix

\newcommand{\Stinespring}{\br{\pi, V, \dil}}

% ---------------------------------------------------------------------------
% Title and authors
% ---------------------------------------------------------------------------
\title{\textbf{A New Proof of a Generalized Stinespring Dilation Theorem}%
\thanks{Acknowledgements common to all authors.}}

\author{
  Author One\thanks{Affiliation details, address, email.} \\
  Affiliation
  \and
  Author Two\thanks{Affiliation details, address, email.} \\
  U.\ Waterloo
}

\date{\today}

\begin{document}

\maketitle

\begin{abstract}
The Stinespring dilation theorem~\cite{Stinespring1955} represents every
completely positive map on a C\(^*\)-algebra as the compression of a
\(*\)-representation. For a finite-dimensional C\(^*\)-algebra \(\alg\) with a
faithful tracial state \(\tau\), we realize the dilation of a linear map
\(\varphi : \alg \to \bdd{\hilb}\) inside the fixed Hilbert space
\(\hilb \tensor L^2(\alg, \tau)\). Stinespring's \(\varphi\)-inner product is
implemented by a single operator \(\Fphi\). It lies in the commutant of the left
regular representation, so in \(\bdd{\hilb} \tensor \alg\op\). It is positive
exactly when \(\varphi\) is completely positive, and in a matrix-unit basis it is
the Choi matrix of \(\varphi\). This replaces the quotient-and-completion step of
the standard proof by the square root \(\Fphi^{1/2}\), and recovers Choi's
theorem as a corollary. \todo{one sentence on the generalization beyond finite
dimensions, if any.}
\end{abstract}

\tableofcontents

% ===========================================================================
\section{Introduction}
\label{sec:introduction}

\paragraph{Background.}
Stinespring's theorem~\cite{Stinespring1955} writes every completely positive
map \(\varphi : \alg \to \bdd{\hilb}\) on a unital C\(^*\)-algebra as
\(\varphi(a) = V^\adjoint \pi(a) V\), with \(\pi\) a \(*\)-representation of
\(\alg\) on a larger Hilbert space and \(V\) bounded. It has many refinements:
minimality and uniqueness up to unitary equivalence~\cite{Arveson1969}, the
completely bounded version, and the Hilbert C\(^*\)-module version of
Kasparov~\cite{Kasparov1980}. A standard reference is \cite{Paulsen2002}. In finite dimensions the closely related result is Choi's
theorem~\cite{Choi1975}: a map on matrices is completely positive if and only if
its Choi matrix is positive semidefinite.

\paragraph{This work.}
Stinespring's proof builds the dilation space abstractly, as the completion of
\(\alg \algtensor \hilb\) modulo the null space of a \(\varphi\)-dependent
sesquilinear form. When \(\alg\) is finite dimensional we fix instead the
concrete space \(\Ltn\), where \(\Ltwo\) is the GNS space of a faithful tracial
state as in the theory of finite von Neumann algebras~\cite{AnantharamanPopa}.
All the dependence on \(\varphi\) is then carried by one positive operator
\(\Fphi\) on this space. We show that \(\Fphi\) commutes with the left regular
representation, so it lies in \(\bdd{\hilb} \tensor \alg\op\), and that it is the
Choi matrix of \(\varphi\).

\paragraph{Contributions.}
\begin{enumerate}[label=(\arabic*)]
  \item A Stinespring dilation of any completely positive
    \(\varphi : \alg \to \bdd{\hilb}\), for \(\alg\) finite dimensional, realized
    inside \(\Ltn\) (\cref{thm:main}).
  \item The identification of the implementing operator \(\Fphi\) with the Choi
    matrix (\cref{thm:choi}), giving a proof of Choi's theorem
    (\cref{cor:choi}).
  \item An extension to arbitrary linear maps (\cref{sec:linear}). The map
    \(\varphi \mapsto \Fphi\) is a linear bijection onto the commutant
    \(\bdd{\hilb} \tensor \alg\op\), and every linear \(\varphi\) satisfies
    \(\varphi(a) = V^\adjoint\rho(a)\Fphi V\) (\cref{thm:linear}). This gives a
    Jordan decomposition and a Krein-space dilation of Hermitian maps
    (\cref{thm:jordan,cor:krein}), and a two-sided dilation of general ones
    (\cref{prop:two-sided}).
  \item \todo{Generalization beyond the finite-dimensional case, e.g.\ to finite
    von Neumann algebras (\cref{sec:open}).}
\end{enumerate}

\paragraph{Organization.}
\Cref{sec:preliminaries} fixes notation and recalls the classical theorem.
\Cref{sec:standard-form} sets up \(\Ltwo\) and the operators acting on it.
\Cref{sec:main-results} states the results, and \cref{sec:proof} proves them.
\Cref{sec:choi} identifies \(\Fphi\) with the Choi matrix, and
\cref{sec:linear} extends the construction to arbitrary linear maps.
\Cref{sec:comparison} compares the construction with Stinespring's.
\Cref{sec:applications} and \cref{sec:open} discuss applications and open
questions.

\paragraph{Acknowledgements.}
Personal acknowledgements, grants, etc.

% ===========================================================================
\section{Preliminaries}
\label{sec:preliminaries}

\subsection{Notation}
\label{sec:notation}

Hilbert spaces are complex and inner products \(\ip{\cdot}{\cdot}\) are linear in
the \emph{first} argument. For vectors \(u, v\) we write \(\rankone{u}{v}\) for the
rank-one operator \(w \mapsto \ip{w}{v} u\). \(\bdd{\hilb}\) denotes the bounded
operators on \(\hilb\), and \(\matr{n}{\alg}\) the \(n \times n\) matrices over
\(\alg\), with matrix units \(E_{ij}\). For a linear map
\(\varphi : \alg \to \algB\) write
\(\varphi_n \eqdef \id_{\mathrm{M}_n} \tensor \varphi :
\matr{n}{\alg} \to \matr{n}{\algB}\).
For a set \(S\) of operators, \(S'\) denotes its commutant.
The opposite algebra \(\alg\op\) is \(\alg\) with product
\(a\op b\op \eqdef (ba)\op\).

\begin{definition}[Complete positivity]
\label{def:cp}
A linear map \(\varphi : \alg \to \algB\) between C\(^*\)-algebras is
\emph{completely positive} if \(\varphi_n\) is positive for every \(n \ge 1\).
\end{definition}

\begin{definition}[Stinespring dilation]
\label{def:dilation}
Let \(\varphi : \alg \to \bdd{\hilb}\) be linear. A \emph{Stinespring dilation} of
\(\varphi\) is a triple \(\Stinespring\) consisting of a Hilbert space \(\dil\), a
unital \(*\)-representation \(\pi : \alg \to \bdd{\dil}\) and a bounded operator
\(V : \hilb \to \dil\) such that
\begin{equation}
\label{eq:dilation}
  \varphi(a) = V^\adjoint \pi(a) V \qquad \text{for all } a \in \alg .
\end{equation}
It is \emph{minimal} if \(\Span\set{\pi(a) V h : a \in \alg,\ h \in \hilb}\) is
dense in \(\dil\).
\end{definition}

\subsection{The classical theorem}
\label{sec:classical}

\begin{theorem}[Stinespring~{\cite[Theorem~1]{Stinespring1955}}]
\label{thm:stinespring}
Let \(\alg\) be a unital C\(^*\)-algebra and \(\varphi : \alg \to \bdd{\hilb}\)
completely positive. Then there are a Hilbert space \(\dil\), a
\(*\)-representation \(\pi : \alg \to \bdd{\dil}\) and a bounded operator
\(V : \hilb \to \dil\) such that \(\varphi(a) = V^\adjoint \pi(a) V\) for all
\(a \in \alg\), and \(\norm{V}^2 = \norm{\varphi(1)}\).
\end{theorem}
\todo{Check the theorem number and the exact statement against the original
paper.}

\begin{remark}[Minimality and uniqueness]
\label{rem:uniqueness}
Replacing \(\dil\) by the closed span of \(\pi(\alg)V\hilb\) gives a minimal
dilation. Any two minimal dilations are unitarily equivalent: there is a unitary
\(U : \dil_1 \to \dil_2\) with \(U V_1 = V_2\) and
\(U \pi_1(a) = \pi_2(a) U\) for all \(a\). These refinements are not in
Stinespring's paper; see~\cite{Arveson1969} and~\cite[Chapter~4]{Paulsen2002}.
\end{remark}

\begin{remark}[The standard construction]
\label{rem:standard}
Stinespring's proof of \cref{thm:stinespring}~\cite{Stinespring1955} equips
\(\hilb \algtensor \alg\) with the sesquilinear form
\begin{equation}
\label{eq:phi-form}
  \ipphi{x \tensor a}{y \tensor b} \eqdef \ip{\varphi(b^\adjoint a) x}{y}_{\hilb},
\end{equation}
uses complete positivity to show that it is positive semidefinite, and takes the
quotient by its null space followed by the completion. Then \(\pi\) acts by left
multiplication on the \(\alg\) factor and \(V x \eqdef [x \tensor 1]\).
\end{remark}

% ===========================================================================
\section{The Standard Form of a Finite-Dimensional C\texorpdfstring{\(^*\)}{*}-Algebra}
\label{sec:standard-form}

From now on \(\alg\) is a finite-dimensional C\(^*\)-algebra.

\begin{theorem}[Structure of finite-dimensional C\(^*\)-algebras]
\label{thm:structure}
There are integers \(n_1, \dots, n_r \ge 1\) and a \(*\)-isomorphism
\[
  \alg \cong \bigoplus_{k=1}^r \mathrm{M}_{n_k}(\complex).
\]
\end{theorem}
See~\cite[Theorem~III.1.1]{Davidson1996} or~\cite{Takesaki1979}.
\todo{Check the theorem number in Davidson and add a precise reference in
Takesaki, Vol.~I.}

We write \(a = (a_1, \dots, a_r)\) for elements of \(\alg\). Let
\(\iota_k : \mathrm{M}_{n_k}(\complex) \to \alg\) be the inclusion as the \(k\)-th
summand and \(p_k : \alg \to \mathrm{M}_{n_k}(\complex)\) the projection onto it,
so that \(a = \sum_k \iota_k(p_k(a))\). For a linear map
\(\varphi : \alg \to \bdd{\hilb}\) set
\[
  \varphi_k \eqdef \varphi \compose \iota_k : \mathrm{M}_{n_k}(\complex) \to \bdd{\hilb}.
\]

\begin{lemma}[Reduction to matrix factors]
\label{lem:reduction}
Let \(\varphi : \alg \to \bdd{\hilb}\) be linear. Then
\begin{enumerate}[label=(\roman*)]
  \item \(\varphi = \sum_{k=1}^r \varphi_k \compose p_k\), so \(\varphi\) is
    determined by \(\varphi_1, \dots, \varphi_r\);
  \item \(\varphi\) is completely positive if and only if every \(\varphi_k\) is.
\end{enumerate}
\end{lemma}
\begin{proof}
(i) is linearity of \(\varphi\). For (ii), \(\iota_k\) and \(p_k\) are
\(*\)-homomorphisms (\(\iota_k\) is not unital). So their amplifications
\((\iota_k)_n\) and \((p_k)_n\) are \(*\)-homomorphisms, hence positive, and
\(\iota_k\), \(p_k\) are completely positive. If \(\varphi\) is completely positive
then so is \(\varphi_k = \varphi \compose \iota_k\). Conversely, if every
\(\varphi_k\) is completely positive then so is the sum
\(\sum_k \varphi_k \compose p_k = \varphi\).
\end{proof}

By \cref{lem:reduction}, statements about a single map on \(\alg\) reduce to
statements about maps on one matrix factor \(\mathrm{M}_n(\complex)\). There we
write \(E_{ij}\) for the matrix units and \(e_i\) for the standard basis of
\(\complex^n\).

\begin{lemma}[Faithful tracial states]
\label{lem:trace}
For weights \(t_1, \dots, t_r > 0\) with \(\sum_k t_k n_k = 1\), the functional
\[
  \tau(a) \eqdef \sum_{k=1}^r t_k \trace(a_k)
\]
is a faithful tracial state on \(\alg\), and every faithful tracial state has this
form.
\end{lemma}
\begin{proof}
Positivity, \(\tau(1) = \sum_k t_k n_k = 1\) and \(\tau(ab) = \tau(ba)\) follow
from the corresponding properties of \(\trace\) on each summand. If
\(\tau(a^\adjoint a) = 0\) then \(\trace(a_k^\adjoint a_k) = 0\) for every \(k\),
so \(a = 0\). Conversely, a tracial functional on \(\mathrm{M}_{n}(\complex)\) is
a multiple of \(\trace\), so any tracial state is of this form with
\(t_k \ge 0\), and faithfulness forces \(t_k > 0\).
\end{proof}

\begin{remark}
The faithful tracial state is unique exactly when \(\alg\) is a factor
(\(r = 1\)), in which case \(\tau = \frac{1}{n}\trace\). Natural choices in general
are \(t_k = 1/(r n_k)\), or \(t_k = n_k / \sum_j n_j^2\), which gives
\(\tau(a) = \trace(\lambda(a)) / \dim \alg\) for the left regular representation
\(\lambda\) below. None of the results below depend on the choice.
\end{remark}

Fix a faithful tracial state \(\tau\) for the rest of the paper.

\begin{definition}[The space \(\Ltwo\)]
\label{def:ltwo}
Let \(\Ltwo = L^2(\alg, \tau)\) be the GNS space of \(\tau\), as
in~\cite{AnantharamanPopa}: it is \(\alg\) equipped with the inner product
\[
  \ip{\vect{a}}{\vect{b}}_{\tau} \eqdef \tau(b^\adjoint a),
\]
where \(\vect{a}\) denotes \(a\) viewed as a vector in \(\Ltwo\). Write
\(\norm{\vect{a}}_2 \eqdef \tau(a^\adjoint a)^{1/2}\) and
\(\Omega \eqdef \vect{1}\).
\end{definition}

\begin{lemma}
\label{lem:ltwo}
The map \(a \mapsto \vect{a}\) is a linear bijection \(\alg \to \Ltwo\). No quotient
or completion is needed, since \(\tau\) is faithful and \(\alg\) is finite
dimensional. Moreover
\(\norm{\vect{a}}_2 \le \norm{a}\) for all \(a \in \alg\), and the two norms are
equivalent but in general not equal.
\end{lemma}
\begin{proof}
Faithfulness means the null space \(\set{a : \tau(a^\adjoint a) = 0}\) is zero, so
the GNS quotient is trivial. A finite-dimensional inner product space is
complete. Finally \(\tau(a^\adjoint a) \le \norm{a^\adjoint a}\,\tau(1) =
\norm{a}^2\), and any two norms on a finite-dimensional space are equivalent.
For \(\alg = \mathrm{M}_n(\complex)\) and \(a = E_{11}\) we get
\(\norm{\vect{a}}_2 = n^{-1/2} < 1 = \norm{a}\) when \(n \ge 2\).
\end{proof}

\begin{remark}
So \(\Ltwo\) and \(\alg\) are the same vector space, but the Hilbert space norm
\(\norm{\cdot}_2\) is not the C\(^*\)-norm. The C\(^*\)-norm is not induced by any
inner product unless \(\dim \alg = 1\). What is true is that \(a \mapsto \vect{a}\)
identifies \(\alg\) with \(\Ltwo\) as vector spaces, and we use this
identification throughout.
\end{remark}

\begin{definition}[Left and right regular representations]
\label{def:regular}
For \(a, b \in \alg\) define \(\lambda(a), R_b \in \bdd{\Ltwo}\) by
\[
  \lambda(a)\vect{c} \eqdef \vect{ac}, \qquad R_b \vect{c} \eqdef \vect{cb}.
\]
\end{definition}

\begin{lemma}
\label{lem:regular}
\begin{enumerate}[label=(\roman*)]
  \item \(\lambda : \alg \to \bdd{\Ltwo}\) is a faithful unital
    \(*\)-representation, and \(\Omega\) is cyclic:
    \(\lambda(\alg)\Omega = \Ltwo\).
  \item \(b\op \mapsto R_b\) is a faithful unital \(*\)-representation of
    \(\alg\op\) on \(\Ltwo\).
  \item \(\lambda(\alg)' = R(\alg\op)\) and \(R(\alg\op)' = \lambda(\alg)\).
\end{enumerate}
\end{lemma}
\begin{proof}
(i) Multiplicativity and unitality are clear, and \(\lambda(a)\Omega = \vect{a}\)
gives cyclicity and faithfulness. For adjoints,
\(\ip{\lambda(c)\vect{a}}{\vect{b}}_{\tau} = \tau(b^\adjoint c a) =
\tau((c^\adjoint b)^\adjoint a) = \ip{\vect{a}}{\lambda(c^\adjoint)\vect{b}}_{\tau}\).

(ii) \(R_b R_c \vect{d} = \vect{dcb} = R_{cb}\vect{d}\), and
\(R(b\op c\op) = R((cb)\op) = R_{cb}\), so \(R\) is multiplicative on \(\alg\op\).
For adjoints we use that \(\tau\) is tracial:
\(\ip{R_b\vect{c}}{\vect{d}}_{\tau} = \tau(d^\adjoint c b) =
\tau(b d^\adjoint c) = \tau((d b^\adjoint)^\adjoint c) =
\ip{\vect{c}}{R_{b^\adjoint}\vect{d}}_{\tau}\).
Faithfulness follows from \(R_b\Omega = \vect{b}\).

(iii) \(\lambda(a)\) and \(R_b\) commute by associativity, so
\(R(\alg\op) \subseteq \lambda(\alg)'\). Conversely let
\(T \in \lambda(\alg)'\), and by \cref{lem:ltwo} let \(d \in \alg\) be such that
\(T\Omega = \vect{d}\). Then for all \(a\),
\[
  T\vect{a} = T\lambda(a)\Omega = \lambda(a)T\Omega = \vect{ad} = R_d\vect{a},
\]
so \(T = R_d\). The second equality follows by the same argument with left and
right exchanged.
\end{proof}

\begin{lemma}[Block decomposition of \(\Ltwo\)]
\label{lem:ltwo-blocks}
Let \(\tau = \sum_k t_k \trace\) be as in \cref{lem:trace}, let
\(\tau_k \eqdef \frac{1}{n_k}\trace\) be the tracial state of
\(\mathrm{M}_{n_k}(\complex)\), and let \(c_k \eqdef \tau(\iota_k(1)) = t_k n_k\).
Then the map
\[
  U : \Ltwo \to \bigoplus_{k=1}^r L^2\br{\mathrm{M}_{n_k}(\complex), \tau_k},
  \qquad \vect{a} \mapsto \br{c_k^{1/2}\,\vect{a_k}}_{k}
\]
is unitary, and for all \(a, b \in \alg\)
\[
  U \lambda(a) U^\adjoint = \bigoplus_k \lambda(a_k), \qquad
  U R_b U^\adjoint = \bigoplus_k R_{b_k}, \qquad
  U\Omega = \br{c_k^{1/2}\,\Omega_k}_k,
\]
where \(\lambda\), \(R\) and \(\Omega_k\) on the right are those of
\(L^2(\mathrm{M}_{n_k}(\complex), \tau_k)\).
\end{lemma}
\begin{proof}
\(\ip{\vect{a}}{\vect{b}}_{\tau} = \sum_k t_k \trace(b_k^\adjoint a_k) =
\sum_k c_k\, \tau_k(b_k^\adjoint a_k)\), so \(U\) is isometric. It is bijective by
\cref{lem:ltwo}. Multiplication in \(\alg\) is componentwise, which gives the
intertwining relations.
\end{proof}

\begin{remark}[\(\bdd{\Ltwo}\) versus \(\alg \tensor \alg\op\)]
\label{rem:not-tensor}
In general \(\bdd{\Ltwo} \ne \lambda(\alg) R(\alg\op) \cong \alg \tensor \alg\op\).
The algebra generated by \(\lambda(\alg)\) and \(R(\alg\op)\) has commutant
\(\lambda(\alg)' \cap R(\alg\op)' = R(\alg\op) \cap \lambda(\alg)\), which is the
centre of \(\alg\). For \(\alg = \complex \oplus \complex\),
\(\alg \tensor \alg\op \cong \complex^4\) is commutative but
\(\bdd{\Ltwo} \cong \mathrm{M}_2(\complex)\). Equality
\(\bdd{\Ltwo} \cong \alg \tensor \alg\op\) holds exactly when \(\alg\) is a factor.
The argument below does not need it: \cref{lem:regular}(iii) suffices.
\end{remark}

\begin{lemma}[Amplified commutant]
\label{lem:amplified-commutant}
For any Hilbert space \(\hilb\),
\[
  \br{\id_{\hilb} \tensor \lambda(\alg)}' = \bdd{\hilb} \tensor R(\alg\op)
  \subseteq \bdd{\Ltn},
\]
where the tensor product on the right is algebraic, which is already closed
because \(R(\alg\op)\) is finite dimensional.
\end{lemma}
\begin{proof}
Since \(\Ltwo\) is finite dimensional, every \(T \in \bdd{\Ltn}\) can be written as
a finite sum \(T = \sum_k S_k \tensor T_k\) with \(S_k \in \bdd{\hilb}\) linearly
independent and \(T_k \in \bdd{\Ltwo}\). (Fix an orthonormal basis of \(\Ltwo\) and
view \(T\) as a matrix over \(\bdd{\hilb}\).) If \(T\) commutes with
\(\id \tensor \lambda(c)\) then \(\sum_k S_k \tensor [T_k, \lambda(c)] = 0\). Taking
functionals \(f_l\) on \(\bdd{\hilb}\) with \(f_l(S_k) = \delta_{kl}\) and slicing
by \(f_l \tensor \id\) gives \([T_l, \lambda(c)] = 0\) for every \(l\). Hence
\(T_l \in R(\alg\op)\) by \cref{lem:regular}(iii). The reverse inclusion is clear.
\end{proof}

% ===========================================================================
\section{Main Results}
\label{sec:main-results}

\subsection{Setting}
\label{sec:setting}

Let \(\alg\) be finite dimensional with faithful tracial state \(\tau\), let
\(\hilb\) be a Hilbert space, and let \(\varphi : \alg \to \bdd{\hilb}\) be linear.
Because \(\Ltwo\) is finite dimensional, \(\Ltn = \hilb \algtensor \Ltwo\), and
the form~\eqref{eq:phi-form} is defined on all of \(\Ltn\):
\begin{equation}
\label{eq:phi-form-ltwo}
  \ipphi{x \tensor \vect{a}}{y \tensor \vect{b}}
  \eqdef \ip{\varphi(b^\adjoint a) x}{y}_{\hilb},
\end{equation}
extended sesquilinearly. Let \(\rho : \alg \to \bdd{\Ltn}\) be the representation
\[
  \rho(a) \eqdef \id_{\hilb} \tensor \lambda(a),
\]
and let \(V : \hilb \to \Ltn\) be
\[
  V x \eqdef x \tensor \Omega .
\]

\begin{remark}
\label{rem:v-alone}
\(V\) is an isometry for the canonical inner product on \(\Ltn\), since
\(\norm{x \tensor \Omega}^2 = \norm{x}^2 \tau(1)\). But
\(V^\adjoint \rho(a) V = \tau(a)\,\id_{\hilb}\), so \((\rho, V, \Ltn)\) dilates the
map \(a \mapsto \tau(a)\id_{\hilb}\), not \(\varphi\). The map \(\varphi\) enters
only through the form~\eqref{eq:phi-form-ltwo}, and the point of the construction
is to move that form into an operator \(\Fphi\).
\end{remark}

\subsection{Statements}
\label{sec:statements}

\begin{theorem}[Stinespring dilation in \(\Ltn\)]
\label{thm:main}
In the setting of \cref{sec:setting}:
\begin{enumerate}[label=(\alph*)]
  \item \label{it:F-exists}
    There is a unique \(\Fphi \in \bdd{\Ltn}\) with
    \(\ipphi{u}{v} = \ip{\Fphi u}{v}\) for all \(u, v \in \Ltn\).
  \item \label{it:F-commutes}
    \(\Fphi\) commutes with \(\rho(\alg)\). Equivalently,
    \(\Fphi \in \bdd{\hilb} \tensor R(\alg\op) \cong \bdd{\hilb} \tensor \alg\op\).
  \item \label{it:F-positive}
    \(\varphi\) is completely positive if and only if \(\Fphi \ge 0\).
  \item \label{it:F-dilation}
    If \(\varphi\) is completely positive, let
    \(\dil \eqdef \overline{\operatorname{ran}}\,\Fphi^{1/2}\),
    \(\pi(a) \eqdef \rho(a)|_{\dil}\) and \(W \eqdef \Fphi^{1/2} V\). Then
    \((\pi, W, \dil)\) is a minimal Stinespring dilation of \(\varphi\), with
    \(W^\adjoint W = \varphi(1)\). In particular \(\norm{W}^2 = \norm{\varphi(1)}\),
    and \(W\) is an isometry if and only if \(\varphi\) is unital.
\end{enumerate}
\end{theorem}

\begin{theorem}[\(\Fphi\) is the Choi matrix]
\label{thm:choi}
Let \(\alg = \mathrm{M}_n(\complex)\) with \(\tau = \frac{1}{n}\trace\), and let
\(\varphi : \mathrm{M}_n(\complex) \to \bdd{\hilb}\) be linear, with Choi matrix
\[
  \Choi \eqdef \sum_{j,l=1}^{n} E_{jl} \tensor \varphi(E_{jl})
  \in \mathrm{M}_n(\complex) \tensor \bdd{\hilb}.
\]
Under the unitary identification \(\Ltwo \cong \complex^n \tensor \complex^n\),
\(\vect{E_{ij}} \mapsto n^{-1/2}\, e_i \tensor e_j\), and after moving \(\hilb\)
to the last tensor factor,
\[
  \Fphi = n\,\br{\id_n \tensor \Choi}.
\]
\end{theorem}

\begin{corollary}[Direct sums]
\label{cor:choi-blocks}
Let \(\alg = \bigoplus_k \mathrm{M}_{n_k}(\complex)\) with \(\tau\) as in
\cref{lem:trace}. Under the unitary of \cref{lem:ltwo-blocks} followed by
\cref{thm:choi} on each summand,
\[
  \Fphi = \bigoplus_{k=1}^r t_k^{-1}\br{\id_{n_k} \tensor C_{\varphi_k}}.
\]
\end{corollary}

\begin{corollary}[Choi~\cite{Choi1975}]
\label{cor:choi}
A linear map \(\varphi : \mathrm{M}_n(\complex) \to \bdd{\hilb}\) is completely
positive if and only if \(\Choi \ge 0\).
\end{corollary}
\begin{proof}
By \cref{thm:choi}, \(\Fphi \ge 0\) if and only if \(\Choi \ge 0\). Apply
\cref{thm:main}\ref{it:F-positive}.
\end{proof}

By \cref{lem:reduction}(ii) and \cref{cor:choi-blocks}, the same holds for
\(\varphi\) on a direct sum: \(\varphi\) is completely positive if and only if
every \(C_{\varphi_k}\) is positive.

% ===========================================================================
\section{Proof of the Main Theorem}
\label{sec:proof}

\subsection{Proof strategy}
\label{sec:strategy}

Stinespring's dilation space is the separated completion of \(\hilb \algtensor \alg\)
in the \(\varphi\)-form. In finite dimensions \(\hilb \algtensor \alg = \Ltn\) is
already a Hilbert space for its canonical inner product, and the \(\varphi\)-form is
a bounded form on it. So it is implemented by an operator \(\Fphi\) and the
separated completion becomes \(\overline{\operatorname{ran}}\,\Fphi^{1/2}\).
Everything Stinespring needs then follows from two facts. The first is that
\(\Fphi \ge 0\), which is Stinespring's positivity argument. The second is that
\(\Fphi\) commutes with left multiplication, which follows from
\(\ipphi{\rho(c) u}{v} = \ipphi{u}{\rho(c^\adjoint) v}\).

Throughout, fix an orthonormal basis \(\set{\vect{u}_p}_{p=1}^m\) of \(\Ltwo\), so
\(u_p \in \alg\), \(\tau(u_q^\adjoint u_p) = \delta_{pq}\) and \(m = \dim \alg\).
Every \(u \in \Ltn\) is uniquely \(u = \sum_p x_p \tensor \vect{u}_p\) with
\(\norm{u}^2 = \sum_p \norm{x_p}^2\).

\subsection{Step 1: The implementing operator}
\label{sec:step-F}

\begin{lemma}
\label{lem:form-bounded}
The form~\eqref{eq:phi-form-ltwo} is well defined and bounded on \(\Ltn\).
\end{lemma}
\begin{proof}
The map \((x, a, y, b) \mapsto \ip{\varphi(b^\adjoint a)x}{y}\) is linear in
\(x, a\) and conjugate linear in \(y, b\). So it induces a unique sesquilinear form
on \(\hilb \algtensor \alg = \Ltn\). In the basis above, for
\(u = \sum_p x_p \tensor \vect{u}_p\) and \(v = \sum_q y_q \tensor \vect{u}_q\),
\[
  \abs{\ipphi{u}{v}}
  = \abs{\sum_{p,q} \ip{\varphi(u_q^\adjoint u_p) x_p}{y_q}}
  \le M \sum_p \norm{x_p} \sum_q \norm{y_q}
  \le M m \norm{u}\norm{v},
\]
with \(M \eqdef \max_{p,q}\norm{\varphi(u_q^\adjoint u_p)}\).
\end{proof}

\begin{proofof}{\cref{thm:main}\ref{it:F-exists}}
A bounded sesquilinear form on a Hilbert space is implemented by a unique bounded
operator~\cite[Theorem~II.2.2]{Conway1990}. Apply this to
\cref{lem:form-bounded}.
\todo{Check the theorem number in Conway.}
\end{proofof}

\subsection{Step 2: Commutation with the left regular representation}
\label{sec:step-commute}

\begin{proofof}{\cref{thm:main}\ref{it:F-commutes}}
For simple tensors and \(c \in \alg\),
\begin{align*}
  \ip{\Fphi \rho(c)(x \tensor \vect{a})}{y \tensor \vect{b}}
  &= \ipphi{x \tensor \vect{ca}}{y \tensor \vect{b}}
   = \ip{\varphi(b^\adjoint c a) x}{y} \\
  &= \ipphi{x \tensor \vect{a}}{y \tensor \vect{c^\adjoint b}}
   = \ip{\Fphi (x \tensor \vect{a})}{\rho(c^\adjoint)(y \tensor \vect{b})} \\
  &= \ip{\rho(c)\Fphi (x \tensor \vect{a})}{y \tensor \vect{b}},
\end{align*}
using \(\rho(c^\adjoint) = \rho(c)^\adjoint\) (\cref{lem:regular}(i)). Simple
tensors span \(\Ltn\), so \(\Fphi\rho(c) = \rho(c)\Fphi\). The equivalent form
follows from \cref{lem:amplified-commutant}.
\end{proofof}

\subsection{Step 3: Positivity}
\label{sec:step-positive}

\begin{proofof}{\cref{thm:main}\ref{it:F-positive}}
Suppose \(\varphi\) is completely positive, and let
\(u = \sum_{p=1}^m x_p \tensor \vect{a}_p\). Then
\[
  \ip{\Fphi u}{u} = \ipphi{u}{u} = \sum_{p,q} \ip{\varphi(a_q^\adjoint a_p)x_p}{x_q}
  = \ip{\varphi_m(P)\,\vec{x}}{\vec{x}}_{\hilb^m},
\]
where \(\vec{x} = (x_1, \dots, x_m)\) and \(P = [a_q^\adjoint a_p]_{q,p} =
X^\adjoint X \in \matr{m}{\alg}\) for the row \(X = (a_1, \dots, a_m)\). Since
\(P \ge 0\) and \(\varphi_m\) is positive, \(\ip{\Fphi u}{u} \ge 0\). This is
Stinespring's argument~\cite{Stinespring1955}; see
also~\cite[Theorem~4.1]{Paulsen2002}.

Conversely, if \(\Fphi \ge 0\) then part~\ref{it:F-dilation}, whose proof below
uses only \(\Fphi \ge 0\), writes \(\varphi(a) = W^\adjoint \pi(a) W\) with \(\pi\)
a \(*\)-representation. Such a map is completely positive, because
\(\varphi_n(P) = (\id_n \tensor W)^\adjoint \pi_n(P) (\id_n \tensor W)\) and
\(\pi_n\) is a \(*\)-homomorphism.
\end{proofof}

\subsection{Step 4: The dilation}
\label{sec:step-dilation}

\begin{proofof}{\cref{thm:main}\ref{it:F-dilation}}
Assume \(\Fphi \ge 0\). By~\ref{it:F-commutes} and the continuous functional
calculus, \(\Fphi^{1/2}\) also commutes with \(\rho(\alg)\). Hence
\(\dil = \overline{\operatorname{ran}}\,\Fphi^{1/2}\) is invariant under
\(\rho(\alg)\), which is a self-adjoint set, and \(\pi = \rho|_{\dil}\) is a
unital \(*\)-representation. For \(x, y \in \hilb\) and \(a \in \alg\),
\begin{align*}
  \ip{W^\adjoint \pi(a) W x}{y}
  &= \ip{\rho(a)\Fphi^{1/2}(x \tensor \Omega)}{\Fphi^{1/2}(y \tensor \Omega)}
   = \ip{\Fphi (x \tensor \vect{a})}{y \tensor \Omega} \\
  &= \ipphi{x \tensor \vect{a}}{y \tensor \vect{1}}
   = \ip{\varphi(a)x}{y}.
\end{align*}
So \(\varphi(a) = W^\adjoint \pi(a) W\), and taking \(a = 1\) gives
\(W^\adjoint W = \varphi(1)\). For minimality,
\[
  \Span \pi(\alg) W \hilb = \Fphi^{1/2}\,\Span\set{x \tensor \vect{a}}
  = \Fphi^{1/2}(\Ltn),
\]
which is dense in \(\dil\) by definition.
\end{proofof}

% ===========================================================================
\section{Identification with the Choi Matrix}
\label{sec:choi}

\begin{lemma}[\(\Fphi\) in an orthonormal basis]
\label{lem:F-basis}
With \(\set{\vect{u}_p}\) as in \cref{sec:strategy},
\[
  \Fphi = \sum_{p,q=1}^m \varphi(u_p^\adjoint u_q) \tensor
  \rankone{\vect{u}_p}{\vect{u}_q}.
\]
\end{lemma}
\begin{proof}
Both sides are bounded, so it suffices to compare matrix entries on simple tensors
\(x \tensor \vect{u}_i\), \(y \tensor \vect{u}_j\). The right-hand side gives
\(\sum_{p,q}\ip{\varphi(u_p^\adjoint u_q)x}{y}\,\delta_{qi}\delta_{pj} =
\ip{\varphi(u_j^\adjoint u_i)x}{y}\), which is
\(\ipphi{x \tensor \vect{u}_i}{y \tensor \vect{u}_j}\).
\end{proof}

\subsection{The factor case}
\label{sec:choi-factor}

\begin{proofof}{\cref{thm:choi}}
Since \(\tau(E_{ij}^\adjoint E_{i'j'}) = \frac{1}{n}\delta_{ii'}\delta_{jj'}\),
the elements \(u_{ij} \eqdef n^{1/2}E_{ij}\) form an orthonormal basis
\(\set{\vect{u}_{ij}}\) of \(\Ltwo\), and
\[
  u_{ij}^\adjoint\, u_{i'l} = n\,\delta_{ii'}\,E_{jl}.
\]
Substituting into \cref{lem:F-basis}, only the terms with equal first index
\(i = i'\) survive:
\[
  \Fphi = n \sum_{j,l} \varphi(E_{jl}) \tensor
  \sum_i \rankone{\vect{u}_{ij}}{\vect{u}_{il}}.
\]
Under \(\vect{u}_{ij} \mapsto e_i \tensor e_j\), we have
\(\sum_i \rankone{\vect{u}_{ij}}{\vect{u}_{il}} \mapsto \id_n \tensor E_{jl}\).
Moving \(\hilb\) to the last factor gives \(\Fphi = n(\id_n \tensor \Choi)\).
\end{proofof}

\begin{remark}[Consistency with \cref{thm:main}\ref{it:F-commutes}]
\label{rem:choi-consistency}
Under the identification of \cref{thm:choi}, \(\lambda(a) = a \tensor \id_n\) and
\(R_b = \id_n \tensor b^\transpose\) on \(\complex^n \tensor \complex^n\). So the
factor \(\id_n\) in \(\Fphi\) is the leg on which \(\lambda\) acts. The Choi
matrix sits in the leg \(R(\alg\op) \cong \alg\op\), where \(\alg\op \cong \alg\)
via \(b\op \mapsto b^\transpose\), exactly as
\cref{thm:main}\ref{it:F-commutes} predicts. As an element of
\(\bdd{\hilb} \tensor \alg\op\),
\(\Fphi = n\sum_{j,l}\varphi(E_{jl}) \tensor E_{lj}\op\).
\end{remark}

\subsection{Reduction of direct sums to factors}
\label{sec:choi-blocks}

\begin{proposition}
\label{prop:F-blocks}
Let \(\alg = \bigoplus_k \mathrm{M}_{n_k}(\complex)\), with \(\tau\), \(\tau_k\),
\(c_k\) and \(U\) as in \cref{lem:ltwo-blocks}. For each \(k\), let
\(F_{\varphi_k}\) be the operator of \cref{thm:main}\ref{it:F-exists} for
\(\varphi_k\) on \(\hilb \tensor L^2(\mathrm{M}_{n_k}(\complex), \tau_k)\). Then
\[
  (\id_{\hilb} \tensor U)\,\Fphi\,(\id_{\hilb} \tensor U)^\adjoint
  = \bigoplus_{k=1}^r c_k^{-1}\, F_{\varphi_k}.
\]
\end{proposition}
\begin{proof}
Since \(b^\adjoint a = \sum_k \iota_k(b_k^\adjoint a_k)\), \cref{lem:reduction}(i)
gives
\[
  \ipphi{x \tensor \vect{a}}{y \tensor \vect{b}}
  = \sum_k \ip{\varphi_k(b_k^\adjoint a_k)x}{y}
  = \sum_k \ip{x \tensor \vect{a_k}}{y \tensor \vect{b_k}}_{\varphi_k}.
\]
The \(k\)-th component of \((\id \tensor U)(x \tensor \vect{a})\) is
\(c_k^{1/2}\, x \tensor \vect{a_k}\). So for \(u, v \in \Ltn\), with
\(\tilde{u} = (\id \tensor U)u\) and \(\tilde{v} = (\id \tensor U)v\),
\[
  \ip{\Fphi u}{v} = \ipphi{u}{v}
  = \sum_k c_k^{-1}\ip{\tilde{u}_k}{\tilde{v}_k}_{\varphi_k}
  = \sum_k c_k^{-1}\ip{F_{\varphi_k}\tilde{u}_k}{\tilde{v}_k},
\]
first on simple tensors and then by sesquilinearity. By uniqueness in
\cref{thm:main}\ref{it:F-exists}, this is the claim.
\end{proof}

\begin{proofof}{\cref{cor:choi-blocks}}
By \cref{prop:F-blocks} and \cref{thm:choi} applied to each \(\varphi_k\),
\(\Fphi\) is unitarily equivalent to
\(\bigoplus_k c_k^{-1} n_k (\id_{n_k} \tensor C_{\varphi_k})\), and
\(c_k^{-1} n_k = t_k^{-1}\).
\end{proofof}

\begin{remark}
\Cref{prop:F-blocks} together with \cref{lem:ltwo-blocks} shows that the whole
construction of \cref{thm:main} splits over the summands: \(\rho\), \(\Fphi\),
\(\dil\) and \(\pi\) are direct sums of the corresponding objects for the
\(\varphi_k\), and \(W x = (c_k^{-1/2}\, c_k^{1/2} F_{\varphi_k}^{1/2}(x \tensor
\Omega_k))_k = (W_k x)_k\). So no generality is lost by working with a single
matrix factor.
\end{remark}

\begin{remark}[Dimension count]
For \(\alg = \mathrm{M}_n(\complex)\) and \(\hilb = \complex^d\), the dilation space
of \cref{thm:main} has
\(\dim \dil = \operatorname{rank}\Fphi = n \cdot \operatorname{rank}\Choi\).
This matches the general form of a minimal dilation of a map on
\(\mathrm{M}_n(\complex)\): every unital \(*\)-representation of
\(\mathrm{M}_n(\complex)\) is a multiple of the identity representation, so
\(\dil \cong \complex^n \tensor \complex^{\kappa}\) with
\(\pi(a) = a \tensor \id_\kappa\). Here the multiplicity is
\(\kappa = \operatorname{rank}\Choi\).
\end{remark}

% ===========================================================================
\section{Extension to Arbitrary Linear Maps}
\label{sec:linear}

Complete positivity is used in \cref{thm:main} only in parts~\ref{it:F-positive}
and~\ref{it:F-dilation}. Parts~\ref{it:F-exists} and~\ref{it:F-commutes}, and the
identification with the Choi matrix in \cref{thm:choi,cor:choi-blocks}, hold for
every linear \(\varphi : \alg \to \bdd{\hilb}\). This is because
\cref{lem:form-bounded}, the commutation argument of \cref{sec:step-commute}
and \cref{lem:F-basis} never use positivity. In this section we develop the
construction for arbitrary linear maps. Throughout, \(\rho\), \(V\), \(\Omega\)
and \(R\) are as in \cref{sec:setting}, and
\[
  \mathcal{C} \eqdef \rho(\alg)' = \bdd{\hilb} \tensor R(\alg\op)
\]
is the commutant, a von Neumann algebra by \cref{lem:amplified-commutant}.

\begin{remark}[\(\mathcal{C}\) as \(\alg \tensor \bdd{\hilb}\)]
\label{rem:op-iso}
Since \(R : \alg\op \to R(\alg\op)\) is a \(*\)-isomorphism (\cref{lem:regular}),
\(\mathcal{C} \cong \alg\op \tensor \bdd{\hilb}\). For finite-dimensional \(\alg\)
we also have \(\alg\op \cong \alg\) as C\(^*\)-algebras: under
\cref{thm:structure}, \(b\op \mapsto (b_1^\transpose, \dots, b_r^\transpose)\) is
a \(*\)-isomorphism, because the transpose reverses products and commutes with
adjoints. Hence \(\mathcal{C} \cong \alg \tensor \bdd{\hilb}\). This isomorphism
depends on the choice of matrix units, and \(\alg\op\) is the canonical object.
\end{remark}

\subsection{The correspondence \texorpdfstring{\(\varphi \leftrightarrow \Fphi\)}{phi <-> F}}
\label{sec:linear-bijection}

\begin{theorem}[Representation of linear maps]
\label{thm:linear}
For every linear map \(\varphi : \alg \to \bdd{\hilb}\), the operator \(\Fphi\)
of \cref{thm:main}\ref{it:F-exists} lies in \(\mathcal{C}\) and satisfies
\begin{equation}
\label{eq:linear-dilation}
  \varphi(a) = V^\adjoint \rho(a) \Fphi V = V^\adjoint \Fphi \rho(a) V
  \qquad \text{for all } a \in \alg .
\end{equation}
Moreover the map \(\varphi \mapsto \Fphi\) is a linear bijection from the space
of linear maps \(\alg \to \bdd{\hilb}\) onto \(\mathcal{C}\). Its inverse sends
\(T \in \mathcal{C}\) to \(\varphi_T(a) \eqdef V^\adjoint \rho(a) T V\).
\end{theorem}
\begin{proof}
\(\Fphi \in \mathcal{C}\) is \cref{thm:main}\ref{it:F-commutes}, whose proof
does not use positivity. For \(x, y \in \hilb\),
\[
  \ip{V^\adjoint \rho(a)\Fphi V x}{y}
  = \ip{\Fphi \rho(a)(x \tensor \Omega)}{y \tensor \Omega}
  = \ipphi{x \tensor \vect{a}}{y \tensor \vect{1}}
  = \ip{\varphi(a)x}{y},
\]
which is~\eqref{eq:linear-dilation}; the second form follows from
\(\Fphi \in \mathcal{C}\).

The form~\eqref{eq:phi-form-ltwo} depends linearly on \(\varphi\), so by
uniqueness in \cref{thm:main}\ref{it:F-exists} so does \(\Fphi\).
By~\eqref{eq:linear-dilation}, \(\Fphi = 0\) forces \(\varphi = 0\), so the map
is injective. For surjectivity, let \(T \in \mathcal{C}\) and
\(\psi \eqdef \varphi_T\). Using \(\rho(b)^\adjoint = \rho(b^\adjoint)\) and
\(T \in \mathcal{C}\),
\[
  \ip{F_\psi (x \tensor \vect{a})}{y \tensor \vect{b}}
  = \ip{\psi(b^\adjoint a)x}{y}
  = \ip{\rho(b^\adjoint a) T (x \tensor \Omega)}{y \tensor \Omega}
  = \ip{T (x \tensor \vect{a})}{y \tensor \vect{b}},
\]
so \(F_\psi = T\).
\end{proof}

\begin{remark}
Equation~\eqref{eq:linear-dilation} says that \((\rho, V, \Ltn)\), which by
\cref{rem:v-alone} is the minimal Stinespring dilation of the tracial map
\(a \mapsto \tau(a)\id_{\hilb}\), dilates \emph{every} linear map, provided the
dilation is twisted by an element of the commutant. For completely positive
\(\varphi\) this is the Radon--Nikodym theorem of
Arveson~\cite{Arveson1969} and Belavkin--Staszewski~\cite{BelavkinStaszewski1986}
with respect to the tracial map, as computed for matrix algebras
in~\cite[Theorem~5.1]{Raginsky2003}. \Cref{thm:linear} removes the positivity
assumption. If \(\hilb = \complex^d\), both sides of the bijection have dimension
\(d^2 \dim \alg\).
\end{remark}

\subsection{Hermitian maps}
\label{sec:hermitian}

For a linear map \(\varphi\) define \(\varphi^\sharp(a) \eqdef \varphi(a^\adjoint)^\adjoint\).
Then \(\varphi\) is \emph{Hermitian} (maps self-adjoint elements to self-adjoint
operators) exactly when \(\varphi^\sharp = \varphi\).

\begin{lemma}
\label{lem:F-adjoint}
\(F_{\varphi^\sharp} = \Fphi^\adjoint\). In particular \(\varphi\) is Hermitian if and
only if \(\Fphi\) is self-adjoint.
\end{lemma}
\begin{proof}
For simple tensors,
\[
  \ip{\Fphi^\adjoint (x \tensor \vect{a})}{y \tensor \vect{b}}
  = \overline{\ip{\Fphi (y \tensor \vect{b})}{x \tensor \vect{a}}}
  = \overline{\ip{\varphi(a^\adjoint b) y}{x}}
  = \ip{\varphi(a^\adjoint b)^\adjoint x}{y}
  = \ip{\varphi^\sharp(b^\adjoint a) x}{y}.
\]
The last expression is \(\ip{F_{\varphi^\sharp}(x \tensor \vect{a})}{y \tensor \vect{b}}\).
The second statement follows from injectivity in \cref{thm:linear}.
\end{proof}

\begin{theorem}[Jordan decomposition]
\label{thm:jordan}
Let \(\varphi\) be Hermitian and let \(\Fphi = F_+ - F_-\) be the Jordan
decomposition of the self-adjoint operator \(\Fphi\), with \(F_\pm \ge 0\) and
\(F_+ F_- = 0\). Then \(F_\pm \in \mathcal{C}\), and \(\varphi_\pm \eqdef \varphi_{F_\pm}\) are
completely positive maps with
\[
  \varphi = \varphi_+ - \varphi_-, \qquad F_{\varphi_\pm} = F_\pm .
\]
\end{theorem}
\begin{proof}
\(F_\pm = f_\pm(\Fphi)\) for the continuous functions \(f_\pm(t) = \max(\pm t, 0)\),
so they lie in the von Neumann algebra \(\mathcal{C}\). By \cref{thm:linear},
\(F_{\varphi_\pm} = F_\pm\) and \(\varphi = \varphi_{\Fphi} = \varphi_+ - \varphi_-\).
Since \(F_\pm \ge 0\), \cref{thm:main}\ref{it:F-positive} shows that \(\varphi_\pm\)
are completely positive.
\end{proof}

\begin{corollary}[Krein-space dilation]
\label{cor:krein}
Let \(\varphi\) be Hermitian. Put \(\dil \eqdef \overline{\operatorname{ran}}\,\abs{\Fphi}\),
\(\pi(a) \eqdef \rho(a)|_{\dil}\), \(W \eqdef \abs{\Fphi}^{1/2} V\), and let \(\Sigma\) be the
restriction to \(\dil\) of the sign of \(\Fphi\). Then \(\Sigma\) is a self-adjoint
unitary on \(\dil\) commuting with \(\pi(\alg)\), and
\[
  \varphi(a) = W^\adjoint \pi(a) \Sigma W \qquad \text{for all } a \in \alg .
\]
\end{corollary}
\begin{proof}
\(\abs{\Fphi}\), \(\abs{\Fphi}^{1/2}\) and \(\operatorname{sgn}(\Fphi)\) all lie in
\(\mathcal{C}\), so \(\dil\) reduces \(\rho(\alg)\) and \(\Sigma\) commutes with \(\pi(\alg)\).
On \(\dil\), \(\operatorname{sgn}(\Fphi)\) squares to the identity. Finally
\(\Fphi = \abs{\Fphi}^{1/2}\operatorname{sgn}(\Fphi)\abs{\Fphi}^{1/2}\), and
\eqref{eq:linear-dilation} gives
\(\varphi(a) = V^\adjoint \abs{\Fphi}^{1/2} \rho(a) \operatorname{sgn}(\Fphi)
\abs{\Fphi}^{1/2} V = W^\adjoint \pi(a) \Sigma W\).
\end{proof}

\begin{remark}
\Cref{cor:krein} says that \(\dil\) with the indefinite form
\(\ip{\Sigma\cdot}{\cdot}\) is a Krein space.
(We write \(\Sigma\) for the fundamental symmetry rather than the customary
\(J\), the standard notation for the conjugation
\(\hat a \mapsto \widehat{a^\adjoint}\) on \(L^2(\alg)\).) Since \(\Sigma\) commutes with \(\pi(\alg)\), the
\(\Sigma\)-adjoint \(\Sigma\pi(a)^\adjoint \Sigma\) of \(\pi(a)\) is \(\pi(a^\adjoint)\), so \(\pi\) is
also a \(*\)-representation for the Krein structure. When \(\varphi\) is completely positive,
\(F_- = 0\), \(\Sigma = \id\), and \cref{cor:krein} reduces to
\cref{thm:main}\ref{it:F-dilation}.
\todo{Locate the existing literature on Krein-space Stinespring dilations of
Hermitian maps and compare.}
\end{remark}

\subsection{General linear maps}
\label{sec:general-linear}

\begin{corollary}
\label{cor:four-cp}
Every linear map \(\varphi : \alg \to \bdd{\hilb}\) is a linear combination of four
completely positive maps:
\[
  \varphi = (\varphi_{1,+} - \varphi_{1,-}) + i\,(\varphi_{2,+} - \varphi_{2,-}),
\]
where \(\varphi_1 = \frac12(\varphi + \varphi^\sharp)\) and
\(\varphi_2 = \frac{1}{2i}(\varphi - \varphi^\sharp)\) are Hermitian and
\(\varphi_{j,\pm}\) are given by \cref{thm:jordan}. Correspondingly,
\(\Fphi = \operatorname{Re}\Fphi + i\operatorname{Im}\Fphi\) with
\(F_{\varphi_1} = \operatorname{Re}\Fphi\) and \(F_{\varphi_2} = \operatorname{Im}\Fphi\).
\end{corollary}
\begin{proof}
\(\varphi_1, \varphi_2\) are Hermitian and \(\varphi = \varphi_1 + i\varphi_2\). By
linearity and \cref{lem:F-adjoint},
\(F_{\varphi_1} = \frac12(\Fphi + \Fphi^\adjoint)\) and
\(F_{\varphi_2} = \frac{1}{2i}(\Fphi - \Fphi^\adjoint)\). Apply \cref{thm:jordan} to each.
\end{proof}

\begin{proposition}[Two-sided dilation]
\label{prop:two-sided}
Let \(\varphi\) be linear and \(\Fphi = U\abs{\Fphi}\) the polar decomposition.
Put \(W_1 \eqdef \abs{\Fphi^\adjoint}^{1/2} V\) and \(W_2 \eqdef U\abs{\Fphi}^{1/2} V\).
Then
\[
  \varphi(a) = W_1^\adjoint \rho(a) W_2 \qquad \text{for all } a \in \alg,
\]
and \(\norm{W_1}\norm{W_2} \le \norm{\Fphi}\).
\end{proposition}
\begin{proof}
\(U\), \(\abs{\Fphi}\) and \(\abs{\Fphi^\adjoint}\) lie in the von Neumann algebra
\(\mathcal{C}\), and \(U\abs{\Fphi}^{1/2} = \abs{\Fphi^\adjoint}^{1/2} U\). Hence
\(\Fphi = \abs{\Fphi^\adjoint}^{1/2} U \abs{\Fphi}^{1/2}\), and
by~\eqref{eq:linear-dilation}
\[
  \varphi(a) = V^\adjoint \abs{\Fphi^\adjoint}^{1/2}\rho(a) U \abs{\Fphi}^{1/2} V
  = W_1^\adjoint \rho(a) W_2 .
\]
For the norm bound, \(\norm{W_1} \le \norm{\Fphi}^{1/2}\) and
\(\norm{W_2} \le \norm{\Fphi}^{1/2}\), since \(V\) is an isometry.
\end{proof}

\begin{remark}
\Cref{prop:two-sided} has the shape of the Wittstock / Haagerup--Paulsen
representation \(\varphi(a) = V_1^\adjoint \pi(a) V_2\) of completely bounded
maps~\cite[Theorem~8.4]{Paulsen2002} \todo{check theorem number}. In finite dimensions every linear map is
completely bounded, so there is no contradiction. In that theorem
\(\norm{V_1}\norm{V_2}\) can be made equal to \(\norm{\varphi}_{cb}\), whereas
\cref{prop:two-sided} only gives the bound \(\norm{\Fphi}\).
\todo{Compare \(\norm{\Fphi}\) with \(\norm{\varphi}_{cb}\). By
\cref{thm:choi}, for \(\alg = \mathrm{M}_n(\complex)\),
\(\norm{\Fphi} = n\norm{\Choi}\). Can the polar decomposition be chosen to attain
\(\norm{\varphi}_{cb}\)?}
\end{remark}

\begin{remark}[Positive but not completely positive maps]
For a positive map \(\varphi\) that is not completely positive, \(\Fphi\) is
self-adjoint but not positive, so \(F_- \ne 0\) in \cref{thm:jordan}.
\todo{Characterize positivity, or \(k\)-positivity, of \(\varphi\) in terms of
\(\Fphi\). For \(\alg = \mathrm{M}_n(\complex)\) this should recover
block-positivity of the Choi matrix.}
\end{remark}

% ===========================================================================
\section{Comparison with Stinespring's Construction}
\label{sec:comparison}

\begin{proposition}
\label{prop:comparison}
Let \(\dil_\varphi\) be the space of the standard construction
(\cref{rem:standard}). The map
\([x \tensor a] \mapsto \Fphi^{1/2}(x \tensor \vect{a})\) extends to a unitary
\(\dil_\varphi \to \dil\) that intertwines the two representations and the two
dilation operators.
\end{proposition}
\begin{proof}
\(\norm{\Fphi^{1/2}u}^2 = \ip{\Fphi u}{u} = \ipphi{u}{u}\), so the map is well
defined on the quotient by the null space and is isometric. Its range
\(\Fphi^{1/2}(\Ltn)\) is dense in \(\dil\). Intertwining follows from
\(\Fphi^{1/2}\rho(a) = \rho(a)\Fphi^{1/2}\) and
\([x \tensor 1] \mapsto \Fphi^{1/2}(x \tensor \Omega) = Wx\).
Alternatively, both dilations are minimal, so this also follows from the
uniqueness of minimal dilations (\cref{rem:uniqueness}).
\end{proof}

\todo{Discuss what is gained: the dilation space is a concrete subspace of a fixed
space independent of \(\varphi\); positivity of the form is positivity of one
operator; the commutant structure explains why the Choi matrix appears; compare
with~\cite{Kasparov1980,Arveson1969,Paulsen2002}.}

% ===========================================================================
\section{Applications}
\label{sec:applications}

\todo{Candidates to check: Kraus decompositions from a spectral decomposition of
\(\Fphi\); a Radon--Nikodym theorem for CP maps \(\psi \le \varphi\) in terms of
\(\pi(\alg)'\), which here is a compression of \(\bdd{\hilb} \tensor R(\alg\op)\);
quantum channels and instruments.}

% ===========================================================================
\section{Open Questions}
\label{sec:open}

\begin{enumerate}[label=(Q\arabic*)]
  \item \emph{Finite von Neumann algebras.} For \(\alg\) a finite von Neumann
    algebra with faithful normal tracial state, \(\Ltwo\) is a genuine completion
    and \(\lambda(\alg)' = R(\alg\op)\) is the standard commutation
    theorem~\cite{AnantharamanPopa}. Is the \(\varphi\)-form still implemented by an
    operator affiliated with \(\bdd{\hilb} \bar{\tensor} \alg\op\), and under
    what conditions on \(\varphi\) is it bounded?
  \item \todo{Non-tracial states (modular theory), or C\(^*\)-algebras without
    a faithful trace.}
  \item \todo{Question.}
\end{enumerate}

\bibliographystyle{plain}
\bibliography{references}

% \appendix
% \input{appendix}

\end{document}
